\documentclass[10pt]{article}

\usepackage[T1]{fontenc}
\usepackage[american]{babel}
\usepackage[a4paper,textwidth=32pc,top=12mm,bottom=14mm]{geometry}
\usepackage{natbib}
\bibliographystyle{compling}
\usepackage{mathtools}
\usepackage{amssymb}
\usepackage{graphicx}
\usepackage{xcolor}
\usepackage{tikz}
\usetikzlibrary{arrows.meta,positioning,calc}
\usepackage[hidelinks]{hyperref}

\makeatletter
\setlength{\@fptop}{0pt}
\makeatother

\hypersetup{
  pdftitle={Ten Creative Problem-Solving Running-Example Proposals},
  pdfsubject={Creative exploration and invention under two truth contracts}
}

\newcommand{\TC}{\textsc{TC}}
\newcommand{\LD}{\textsc{LD}}
\newcommand{\Hall}{\textsc{H}}
\newcommand{\SUP}{\textsc{SUP}}

\definecolor{clrH}{RGB}{155,18,18}
\definecolor{clrLD}{RGB}{22,98,48}
\definecolor{clrB}{RGB}{22,52,125}
\definecolor{clrP}{RGB}{92,61,145}
\definecolor{clrO}{RGB}{170,91,24}
\definecolor{clrG}{RGB}{65,65,65}
\definecolor{bgB}{RGB}{245,247,254}

\tikzset{
  input/.style={rectangle,rounded corners=4pt,draw=gray!55,fill=gray!8,
    font=\small\itshape,text width=5.0cm,align=center,inner sep=6pt},
  decision/.style={rectangle,rounded corners=3pt,draw=clrB!55,fill=bgB,
    font=\small,text width=5.2cm,align=center,inner sep=6pt},
  verdict/.style={rectangle,rounded corners=3pt,draw=clrG,fill=gray!8,
    font=\small\bfseries\sffamily,inner xsep=7pt,inner ysep=4pt},
  verdictH/.style={verdict,draw=clrH,text=clrH,fill=clrH!7},
  verdictLD/.style={verdict,draw=clrLD,text=clrLD,fill=clrLD!7},
  arrow/.style={-{Stealth[length=5pt,width=3.5pt]},thick,color=gray!70!black},
  edge/.style={font=\scriptsize\sffamily,fill=white,inner sep=1.5pt}
}

% Each proposal below is an author-constructed specification check adapted
% from the task design of a mapped resource.  None is represented as a sampled
% model output or as an original item from the cited benchmark.

\newcommand{\RenderProposal}{%
\begin{figure}[!t]
\centering
\resizebox{0.86\linewidth}{!}{%
\begin{tikzpicture}[x=1cm,y=1cm]
  \tikzset{
    taskbox/.style={
      decision,text width=5.35cm,minimum height=2.15cm,
      font=\footnotesize,inner sep=5pt},
    contractbox/.style={
      decision,text width=5.35cm,minimum height=1.55cm,
      font=\scriptsize,inner sep=4pt},
    contextbox/.style={
      rectangle,rounded corners=4pt,draw=clrO!60,fill=clrO!6,
      text width=10.8cm,align=left,font=\scriptsize,inner sep=5pt},
    sharedspan/.style={
      input,text width=8.7cm,minimum height=1.25cm,
      font=\footnotesize,inner sep=5pt},
    outcome/.style={
      font=\scriptsize,align=center,text width=3.75cm}
  }

  \draw[dashed,gray!50] (0,0.95) -- (0,-9.55);

  \node[taskbox] (strict) at (-3.0,0)
    {\textbf{\SharedTaskName}\\
     Literal prompt: ``\StrictPrompt''};
  \node[taskbox] (hyp) at (3.0,0)
    {\textbf{\SharedTaskName}\\
     Literal prompt: ``\PermissivePrompt''};

  \node[contextbox] (context) at (0,2.55)
    {\textbf{Task context $p$ --- task oracle $O$}\\[-1pt]
     \ContextText};

  \coordinate (branch) at ($(context.south)+(0,-0.30)$);
  \coordinate (strictAbove) at ($(strict.north)+(0,0.25)$);
  \coordinate (hypAbove) at ($(hyp.north)+(0,0.25)$);

  \draw[semithick,draw=black!55]
    (context.south) -- (branch);

  \draw[arrow]
    (branch) -| (strictAbove)
    -- (strict.north);

  \draw[arrow]
    (branch) -| (hypAbove)
    -- (hyp.north);

  \node[contractbox] (tc1) at (-3.0,-2.85)
    {$\TC(p_1)=(O,\Gamma_1,\mu)$\\
     $O$: \OracleShort\\
     $\Gamma_1=\varnothing$: \StrictScope\\
     $\mu$: uncertainty marked explicitly};
  \node[contractbox] (tc2) at (3.0,-2.85)
    {$\TC(p_2)=(O,\Gamma_2,\mu)$\\
     $O$: \OracleShort\\
     $\Gamma_2$: \PermissiveScope\\
     $\mu$: uncertainty marked explicitly};

  \node[
    rectangle,rounded corners=3pt,draw=clrG!45,fill=gray!4,
    font=\scriptsize,align=center,text width=5.0cm,inner sep=3pt
  ] (sharedstyle) at (0,-4.18)
    {Shared requested style-marking level: \StyleText{}
     (outside the truth contract)};

  \node[sharedspan] (span) at (0,-5.95)
    {\textbf{\upshape Identical response span:}
     ``\ResponseSpan''\\
     \textbf{\upshape Identical claim content:} \ClaimContent;
     \textbf{\upshape observed status marking:} uncertainty marked by
     ``\Marker''};

  \node[verdictH] (hall) at (-3.0,-7.55)
    {HALLUCINATION (H)};
  \node[verdictLD] (licensed) at (3.0,-7.55)
    {LICENSED DIVERGENCE (LD)};
  \node[outcome] at (-3.0,-8.18)
    {\textsc{unknown}; outside $\Gamma_1$\\
     reason code: \textsc{out-of-scope}};
  \node[outcome,text=clrLD] at (3.0,-8.18)
    {\textsc{unknown}; inside $\Gamma_2$\\
     and marked as required};

  \draw[arrow] (strict.south) -- (tc1.north);
  \draw[arrow] (hyp.south) -- (tc2.north);
  \draw[arrow] (tc1.south) -- (span.north -| tc1.south);
  \draw[arrow] (tc2.south) -- (span.north -| tc2.south);
  \draw[arrow] (span.south -| hall.north) -- (hall.north);
  \draw[arrow] (span.south -| licensed.north) -- (licensed.north);

  \node[
    rectangle,rounded corners=3pt,draw=clrG!45,fill=gray!4,
    font=\scriptsize\sffamily\bfseries,align=center,
    text width=7.7cm,inner sep=4pt
  ] at (0,-9.15)
    {Only the permission scope changes the claim label.};
\end{tikzpicture}%
}
\caption{\textbf{Proposal \ProposalNumber: \ProposalTitle.}
The response span, claim content, evidence state, observed status marking,
task oracle, required status marking, and requested style-marking level are
held fixed; only the permission scope changes.  The strict task therefore
assigns the hallucination label (\Hall) with reason code
\textsc{out-of-scope}, whereas the bounded task assigns licensed divergence
(\LD).  The latter records contract compliance, not truth, usefulness,
quality, creativity, or safety.  \SourceNote}
\label{fig:\ProposalLabel}
\end{figure}
}

\begin{document}

% Proposal 1 ---------------------------------------------------------------
\begingroup
\def\ProposalNumber{1}
\def\ProposalTitle{Reaching Cookies in a Deep Jar}
\def\ProposalLabel{creative-cookie-jar}
\def\SharedTaskName{Candidate-solution task}
\def\StrictPrompt{Propose three candidate ways to reach the cookies. Mark every uncertainty explicitly. Do not use untested material properties.}
\def\PermissivePrompt{Propose three candidate ways to reach the cookies. Mark every uncertainty explicitly. You may use one marked, untested material-property hypothesis.}
\def\ContextText{Challenge packet: ``Cookies are at the bottom of a tall, narrow jar. A hand, chopsticks, and tongs cannot reach them. Tape and rubber bands are available.'' No tape-strength test is provided.}
\def\OracleShort{challenge packet and recorded tests}
\def\StrictScope{no untested material properties}
\def\PermissiveScope{marked material hypotheses needed for candidate solutions}
\def\StyleText{concise solution-sketch style}
\def\ResponseSpan{One possibility is that the tape is strong enough to hold a chopstick extension while it lifts one cookie.}
\def\ClaimContent{the tape can hold a chopstick extension while it lifts one cookie}
\def\Marker{One possibility is that}
\def\SourceNote{This author-constructed contrast is adapted from MacGyver's constrained creative-problem-solving setting \citep{tian_macgyver_2025}; the claim and paired contracts are not a benchmark item or sampled model output.}
\RenderProposal
\endgroup
\clearpage

% Proposal 2 ---------------------------------------------------------------
\begingroup
\def\ProposalNumber{2}
\def\ProposalTitle{Searching an Escape Room}
\def\ProposalLabel{creative-escape-room}
\def\SharedTaskName{Search-planning task}
\def\StrictPrompt{Plan the next three searches for the missing lubricant. Mark every location hypothesis explicitly. Do not use unobserved object locations.}
\def\PermissivePrompt{Plan the next three searches for the missing lubricant. Mark every location hypothesis explicitly. You may use marked location hypotheses about visible containers.}
\def\ContextText{Observation at time $t$: ``A rusty key, a locked chest, and a closed desk drawer are visible. The lubricant has not been found, and the drawer has not been searched.'' Hidden simulator state is not admissible evidence under this contract.}
\def\OracleShort{agent observation available at time $t$}
\def\StrictScope{no unobserved object locations}
\def\PermissiveScope{marked search hypotheses about visible containers}
\def\StyleText{concise action-plan style}
\def\ResponseSpan{One possibility is that the lubricant is in the desk drawer.}
\def\ClaimContent{the lubricant is in the desk drawer}
\def\Marker{One possibility is that}
\def\SourceNote{This author-constructed contrast is adapted from EscapeBench's interactive problem-solving setting \citep{qian_escapebench_2025}; its state, claim, and paired contracts are not a benchmark item or sampled model output.}
\RenderProposal
\endgroup
\clearpage

% Proposal 3 ---------------------------------------------------------------
\begingroup
\def\ProposalNumber{3}
\def\ProposalTitle{Designing a Rooftop Community Space}
\def\ProposalLabel{creative-rooftop}
\def\SharedTaskName{Rooftop-concept task}
\def\StrictPrompt{Generate three rooftop concepts. Mark every site assumption explicitly. Do not use unverified site properties.}
\def\PermissivePrompt{Generate three rooftop concepts. Mark every site assumption explicitly. You may use one marked, unverified site-property hypothesis.}
\def\ContextText{Site brief: ``The community center has an unused flat roof.'' The brief contains no load, access, drainage, or sunlight assessment.}
\def\OracleShort{provided community-center site brief}
\def\StrictScope{no unverified site properties}
\def\PermissiveScope{one marked site-property hypothesis for preliminary ideation}
\def\StyleText{plain-language brainstorming style}
\def\ResponseSpan{One possibility is that the roof receives enough summer sunlight to grow herbs.}
\def\ClaimContent{the roof receives enough summer sunlight to grow herbs}
\def\Marker{One possibility is that}
\def\SourceNote{This author-constructed, non-advisory contrast is adapted from CreativityPrism's multi-task creativity setting \citep{hou_creativityprism_2025}; it is not a benchmark item, engineering conclusion, or sampled model output.}
\RenderProposal
\endgroup
\clearpage

% Proposal 4 ---------------------------------------------------------------
\begingroup
\def\ProposalNumber{4}
\def\ProposalTitle{Inventing a Brand Story for a Mug}
\def\ProposalLabel{creative-brand-story}
\def\SharedTaskName{Product-copy task}
\def\StrictPrompt{Write a polished two-sentence blurb about the mug. Mark every undocumented origin explicitly. Do not introduce undocumented product origins.}
\def\PermissivePrompt{Write a polished two-sentence blurb about the mug. Mark every undocumented origin explicitly. You may introduce one marked, invented product-origin hypothesis.}
\def\ContextText{Product sheet: ``Travel mug; blue glazed ceramic; capacity 350 ml.'' It contains no information about the designer, cultural influences, or origin story.}
\def\OracleShort{provided product sheet}
\def\StrictScope{no undocumented product origins}
\def\PermissiveScope{one marked, invented product-origin hypothesis}
\def\StyleText{polished product-copy style}
\def\ResponseSpan{One unverified possibility is that the mug was inspired by Kyoto tea houses.}
\def\ClaimContent{the mug was inspired by Kyoto tea houses}
\def\Marker{One unverified possibility is that}
\def\SourceNote{This author-constructed contrast is adapted from WritingBench's material-conditioned generative-writing setting \citep{wu_writingbench_2025}; it is not a benchmark prompt, product claim, or sampled model output.}
\RenderProposal
\endgroup
\clearpage

% Proposal 5 ---------------------------------------------------------------
\begingroup
\def\ProposalNumber{5}
\def\ProposalTitle{Exploring a Missing-Key Mystery}
\def\ProposalLabel{creative-missing-key}
\def\SharedTaskName{Story-synopsis task}
\def\StrictPrompt{Draft a three-sentence synopsis from the seed. Mark every added backstory explicitly. Do not add events beyond the seed.}
\def\PermissivePrompt{Draft a three-sentence synopsis from the seed. Mark every added backstory explicitly. You may add one marked candidate backstory.}
\def\ContextText{Story seed: ``Mara entered the kitchen. A teapot stood on the counter. The house key was missing.'' The seed gives no account of who moved the key or where it went.}
\def\OracleShort{story seed before any continuation is adopted}
\def\StrictScope{no events beyond the seed}
\def\PermissiveScope{marked candidate backstories for continuation}
\def\StyleText{concise literary-planning style}
\def\ResponseSpan{In one possible continuation, Mara had hidden the missing key inside the teapot.}
\def\ClaimContent{Mara had hidden the missing key inside the teapot}
\def\Marker{In one possible continuation}
\def\SourceNote{This author-constructed contrast is adapted from LitBench's long-form fiction evaluation setting \citep{fein_litbench_2025}; it is evaluated before a continuation becomes canon and is not a benchmark item or sampled model output.}
\RenderProposal
\endgroup
\clearpage

% Proposal 6 ---------------------------------------------------------------
\begingroup
\def\ProposalNumber{6}
\def\ProposalTitle{Explaining a Dark Lighthouse}
\def\ProposalLabel{creative-lighthouse}
\def\SharedTaskName{Five-sentence story-planning task}
\def\StrictPrompt{Plan a five-sentence story from the opening and cue card. Mark every proposed cause explicitly. Do not add causal relations absent from the seed.}
\def\PermissivePrompt{Plan a five-sentence story from the opening and cue card. Mark every proposed cause explicitly. You may add one marked candidate cause.}
\def\ContextText{Opening and cue card: ``At dawn, the lighthouse lamp went dark.'' Cues: lighthouse, violin, storm, letter, window. The card states no cause for the outage.}
\def\OracleShort{opening sentence and cue card}
\def\StrictScope{no new causal relations}
\def\PermissiveScope{marked candidate causes for story development}
\def\StyleText{five-sentence literary style}
\def\ResponseSpan{One possible cause is that the storm cut the power.}
\def\ClaimContent{the storm caused the lighthouse power failure}
\def\Marker{One possible cause is that}
\def\SourceNote{This author-constructed contrast is adapted from the five-sentence, cue-word short-story task studied by \citet{ismayilzada_evaluating_2025}; its seed and paired contracts are not dataset items or sampled model outputs.}
\RenderProposal
\endgroup
\clearpage

% Proposal 7 ---------------------------------------------------------------
\begingroup
\def\ProposalNumber{7}
\def\ProposalTitle{Inventing a Backstory from an Image}
\def\ProposalLabel{creative-station-clock}
\def\SharedTaskName{Literary-caption task}
\def\StrictPrompt{Write a two-sentence literary caption from the image description. Mark every invented backstory explicitly. Do not add events beyond the visible details.}
\def\PermissivePrompt{Write a two-sentence literary caption from the image description. Mark every invented backstory explicitly. You may add one marked invented backstory.}
\def\ContextText{Image description: ``An empty railway platform. A large station clock shows 11:17.'' It states neither that the clock has stopped nor when the last train left.}
\def\OracleShort{provided image description}
\def\StrictScope{no events beyond visible details}
\def\PermissiveScope{one clearly marked invented backstory}
\def\StyleText{evocative literary-caption style}
\def\ResponseSpan{In one possible backstory, the clock stopped at 11:17 when the final train departed.}
\def\ClaimContent{the clock stopped at 11:17 when the final train departed}
\def\Marker{In one possible backstory}
\def\SourceNote{This author-constructed contrast is adapted from the text-to-creative-writing setting of ART or Artifice \citep{chakrabarty_art_2024}; neither the described image nor the paired contracts are a benchmark item or sampled model output.}
\RenderProposal
\endgroup
\clearpage

% Proposal 8 ---------------------------------------------------------------
\begingroup
\def\ProposalNumber{8}
\def\ProposalTitle{Finding New Uses for a Brick}
\def\ProposalLabel{creative-brick}
\def\SharedTaskName{Alternative-use task}
\def\StrictPrompt{List three possible uses for the brick. Mark every untested affordance explicitly. Do not include untested affordances.}
\def\PermissivePrompt{List three possible uses for the brick. Mark every untested affordance explicitly. You may include one marked, untested affordance.}
\def\ContextText{Object card: ``One intact clay brick.'' Supplied tests confirm its use as a doorstop, paperweight, and bookend, but include no smartphone stability test.}
\def\OracleShort{object card and supplied test results}
\def\StrictScope{no untested object affordances}
\def\PermissiveScope{marked hypotheses about alternative uses}
\def\StyleText{brief divergent-thinking list}
\def\ResponseSpan{One possibility is that the brick can support a smartphone at a viewing angle without tipping.}
\def\ClaimContent{the brick can support a smartphone at a viewing angle without tipping}
\def\Marker{One possibility is that}
\def\SourceNote{This author-constructed contrast is adapted from automated divergent-thinking evaluation \citep{organisciak_beyond_2023}; its object card, safety-qualified claim, and paired contracts are not a dataset item or sampled model output.}
\RenderProposal
\endgroup
\clearpage

% Proposal 9 ---------------------------------------------------------------
\begingroup
\def\ProposalNumber{9}
\def\ProposalTitle{Diagnosing a Slow Program}
\def\ProposalLabel{creative-debugging}
\def\SharedTaskName{Timeout-diagnosis task}
\def\StrictPrompt{Write one diagnostic sentence about the timeout. Mark every unverified cause explicitly. Do not introduce causes not established by the log.}
\def\PermissivePrompt{Write one diagnostic sentence about the timeout. Mark every unverified cause explicitly. You may introduce one marked implementation-cause hypothesis.}
\def\ContextText{Execution log: ``The program exceeded the two-second time limit.'' No source code, profiler trace, or operation counts are included.}
\def\OracleShort{supplied execution log}
\def\StrictScope{no causes not established by the log}
\def\PermissiveScope{marked implementation hypotheses for testing}
\def\StyleText{concise technical diagnostic style}
\def\ResponseSpan{One hypothesis is that repeated sorting causes the timeout.}
\def\ClaimContent{repeated sorting causes the timeout}
\def\Marker{One hypothesis is that}
\def\SourceNote{This author-constructed contrast is adapted from NeoCoder's code-generation and execution setting \citep{lu_benchmarking_2025}; its log, diagnosis, and paired contracts are not benchmark items or sampled model outputs.}
\RenderProposal
\endgroup
\clearpage

% Proposal 10 --------------------------------------------------------------
\begingroup
\def\ProposalNumber{10}
\def\ProposalTitle{Searching for a Shorter Geometry Proof}
\def\ProposalLabel{creative-math}
\def\SharedTaskName{Proof-planning task}
\def\StrictPrompt{Write a short note about possible shorter proofs. Mark every unverified strategy explicitly. Do not introduce strategies outside the packet.}
\def\PermissivePrompt{Write a short note about possible shorter proofs. Mark every unverified strategy explicitly. You may introduce one marked conjectured strategy.}
\def\ContextText{Proof packet: one geometry problem and two checked proofs. It contains no reflection construction and no result about whether such a construction yields a shorter valid proof.}
\def\OracleShort{bounded packet of checked proofs}
\def\StrictScope{no proof strategies outside the packet}
\def\PermissiveScope{marked conjectured proof strategies}
\def\StyleText{concise mathematical-planning style}
\def\ResponseSpan{One conjectured route is that a reflection argument yields a shorter valid proof.}
\def\ClaimContent{a reflection argument yields a shorter valid proof}
\def\Marker{One conjectured route is that}
\def\SourceNote{This author-constructed contrast is adapted from CreativeMath's mathematical-creativity setting \citep{ye_assessing_2024}. It treats the supplied proof packet, not complete mathematical truth, as $O$; the example is not a benchmark item or sampled model output.}
\RenderProposal
\endgroup
\clearpage

\bibliography{references}

\end{document}
